% ============================================================
% The Double Transversality
% ============================================================

\documentclass[12pt]{amsart}

\usepackage{amsmath,amssymb,amsthm}
\usepackage{booktabs}
\usepackage{microtype}
\usepackage{url}

\newtheorem{theorem}{Theorem}
\newtheorem{lemma}[theorem]{Lemma}
\newtheorem{proposition}[theorem]{Proposition}
\newtheorem{corollary}[theorem]{Corollary}
\theoremstyle{definition}
\newtheorem{definition}[theorem]{Definition}
\theoremstyle{remark}
\newtheorem{remark}[theorem]{Remark}

\title{The Double Transversality}

\author{Alexander S.\ Petty}
\email{alexander.petty@gmail.com}
\date{October 2025 (revised April 2026)}
\subjclass[2020]{11A63, 11N05, 11M06}

\begin{document}
\begin{abstract}
The preceding paper observed that the collision
invariant's Fourier spectrum is anti-correlated with
the prime character sum magnitudes at lag~$1$: large
collision coefficients tend to pair with characters
having small $|P(s,\chi)|$. This paper reports two
computational extensions of that observation, built on
the exact transform and neutrality framework from the
earlier papers.

First, the anti-correlation persists at every lag. At
lag $\ell$, the collision invariant lives on
$(\mathbb{Z}/b^{\ell+1}\mathbb{Z})^{\times}$ and
involves $\phi(b^{\ell+1})/2$ odd characters. The
Pearson correlation between
$|\hat{S}^{\circ}(\chi)|$ and $|P(s, \chi)|$ is
negative at every base and lag tested, from $3$
characters (base~$3$, lag~$1$) to $2{,}187$ characters
(base~$3$, lag~$7$).

Second, the base sum $F_R(s) = \sum_b w(b)\,
F^{\circ}_b(s)$ nearly vanishes. Individual bases
produce $F^{\circ}_b(0.5)$ values ranging from $-1.4$
to $+5.1$, but their weighted sum is $-0.01$. The
collision invariant's structural content appears to be
base-specific: the individual signals cancel in the
aggregate.

The two observations together suggest a double
transversality: within each base, the Fourier spectrum
is anti-correlated with the zero-sensitive characters;
across bases, the resulting signals point in different
directions and cancel. The term \emph{double
transversality} is used here for that observed
pattern, not for a theorem explaining convergence.
\end{abstract}
\maketitle

% ============================================================
\section{Introduction}

The preceding paper~\cite{paper19} proved the general
neutrality theorem and observed a computational
anti-correlation: the collision coefficients
$|\hat{S}^{\circ}(\chi)|$ are negatively correlated
with the prime character sum magnitudes
$|P(s, \chi)|$ across odd characters at lag~$1$.

This note keeps the status line explicit. The exact
backbone comes from the earlier papers: odd-channel
selection, centering, and the decomposition of
$F^{\circ}_b$ into odd character sums. The new content
here is computational: persistence of the
anti-correlation at higher lags, and near-vanishing of
a finite weighted base sum across a specified set of
prime bases.

Two questions remained. Does the anti-correlation
persist as the lag increases and the family of
$L$-functions grows? And what happens when the
collision invariant is aggregated across bases?

% ============================================================
\section{The Anti-Correlation at Higher Lags}

At lag $\ell$, the collision invariant $S_{\ell}$ lives
on $(\mathbb{Z}/m\mathbb{Z})^{\times}$ with
$m = b^{\ell+1}$. The number of odd characters is
$\phi(m)/2$, which grows as $b^{\ell}(b{-}1)/2$.

\begin{table}[h]
\centering
\begin{tabular}{rrrrrrrr}
\toprule
$b$ & $\ell$ & $m$ & odd & $+$ & $-$ &
net & $\rho$ \\
\midrule
$3$ & $1$ & $9$ & $3$ & $3$ & $0$ &
$+0.62$ & $0.00$ \\
$3$ & $2$ & $27$ & $9$ & $5$ & $4$ &
$+1.18$ & $+0.42$ \\
$3$ & $3$ & $81$ & $27$ & $15$ & $12$ &
$+0.62$ & $-0.42$ \\
$3$ & $4$ & $243$ & $81$ & $35$ & $46$ &
$+4.83$ & $-0.28$ \\
$3$ & $5$ & $729$ & $243$ & $143$ & $100$ &
$+12.8$ & $-0.33$ \\
$3$ & $6$ & $2{,}187$ & $729$ & $359$ & $370$ &
$+29.1$ & $-0.33$ \\
$3$ & $7$ & $6{,}561$ & $2{,}187$ & $1{,}167$ & $1{,}020$ &
$+11.1$ & $-0.20$ \\
\midrule
$5$ & $1$ & $25$ & $10$ & $8$ & $2$ &
$+1.04$ & $-0.45$ \\
$5$ & $2$ & $125$ & $50$ & $34$ & $16$ &
$+4.01$ & $-0.18$ \\
$5$ & $3$ & $625$ & $250$ & $134$ & $116$ &
$+2.48$ & $-0.28$ \\
$5$ & $4$ & $3{,}125$ & $1{,}250$ & $636$ & $614$ &
$+15.3$ & $-0.19$ \\
\midrule
$7$ & $1$ & $49$ & $21$ & $13$ & $8$ &
$+1.81$ & $-0.28$ \\
$7$ & $2$ & $343$ & $147$ & $89$ & $58$ &
$+7.56$ & $-0.22$ \\
$7$ & $3$ & $2{,}401$ & $1{,}029$ & $509$ & $520$ &
$+4.58$ & $-0.16$ \\
\bottomrule
\end{tabular}
\caption{Anti-correlation at higher lags. The Pearson
correlation $\rho$ between $|\hat{S}^{\circ}|$ and
$|P(0.5, \chi)|$ is negative at every lag beyond
$\ell = 2$ at base~$3$. In the tested cases, the net $F^{\circ}(0.5)$ is
positive. $148{,}000$ primes per row.}
\label{tab:lags}
\end{table}

Three patterns emerge.

The anti-correlation persists at every lag. Beyond the
smallest cases (base~$3$, lags~$1$ and~$2$, where
there are too few characters for the correlation to
stabilize), the Pearson coefficient is consistently
negative, ranging from $-0.16$ to $-0.45$.

The net $F^{\circ}(0.5)$ is positive at every base and
lag. From $3$ characters to $2{,}187$ characters, the
sign structure always resolves in favor of convergence.

The net grows with the number of characters. At base~$3$,
it increases from $+0.62$ (lag~$1$) to $+29.1$
(lag~$6$). The convergence does not weaken as the
family of $L$-functions expands.

% ============================================================
\section{The Base Sum}

Let $\mathcal{B} = \{3,5,7,11,13,17,19,23,29,31\}$
be the set of prime bases used in the experiment.
For a fixed prime $p$, the base sum
\[
R(p) = \sum_{b \in \mathcal{B}} w(b)\, S^{\circ}_1(p, b)
\]
aggregates the centered collision deviation across
these bases with weight $w(b) = 1/b^2$. The
base-summed prime harmonic sum is
$F_R(s) = \sum_p R(p) / p^s = \sum_{b \in \mathcal{B}}
w(b)\, F^{\circ}_b(s)$.

\begin{table}[h]
\centering
\begin{tabular}{rrrrr}
\toprule
$s$ & $F_R$ & $\min F^{\circ}_b$ &
$\max F^{\circ}_b$ & range \\
\midrule
$1.0$ & $-0.000$ & $-0.04$ & $+0.21$ & $0.25$ \\
$0.8$ & $-0.001$ & $-0.15$ & $+0.78$ & $0.93$ \\
$0.6$ & $-0.006$ & $-0.66$ & $+2.80$ & $3.46$ \\
$0.5$ & $-0.010$ & $-1.36$ & $+5.07$ & $6.43$ \\
\bottomrule
\end{tabular}
\caption{The base sum $F_R(s)$ compared with the range
of individual $F^{\circ}_b(s)$ across bases
$b = 3, 5, 7, 11, 13, 17, 19, 23, 29, 31$.
$148{,}000$ primes.}
\label{tab:basesum}
\end{table}

At $s = 0.5$, the individual base values span a range
of $6.4$, but their weighted sum is $-0.01$. The
base sum nearly vanishes.

\begin{table}[h]
\centering
\begin{tabular}{rr}
\toprule
$b$ & $F^{\circ}_b(0.5)$ \\
\midrule
$3$ & $-0.03$ \\
$5$ & $-0.00$ \\
$7$ & $+0.07$ \\
$11$ & $-1.36$ \\
$13$ & $-0.18$ \\
$17$ & $-0.09$ \\
$19$ & $-0.55$ \\
$23$ & $-1.30$ \\
$29$ & $+2.80$ \\
$31$ & $+5.07$ \\
\bottomrule
\end{tabular}
\caption{Individual $F^{\circ}_b(0.5)$ values.
Some bases are positive, some negative. Their
weighted sum is $-0.01$.}
\label{tab:ind}
\end{table}

% ============================================================
\section{The Double Transversality}

The computations in the preceding sections suggest two
levels of cancellation.

\emph{Within each base}: the centered sum
$F^{\circ}_b(s)$ is a linear combination of character
sums $P(s, \chi)$ with mixed-sign products. In the
tested cases, the anti-correlation between
$|\hat{S}^{\circ}(\chi)|$ and $|P(s, \chi)|$
accompanies a positive net: the positive terms
outweigh the negative at every base and lag tested.

\emph{Across bases}: the bounded $F^{\circ}_b$ values
have mixed signs across the tested bases, and their
weighted sum over $\mathcal{B}$ nearly vanishes.

\begin{definition}
The \emph{double transversality} is the name for the
observed pattern that:
\begin{enumerate}
\item at each base, the collision spectrum is
anti-correlated with the prime character sum
magnitudes, and
\item across bases, the individual $F^{\circ}_b$
values cancel in the weighted sum.
\end{enumerate}
This is a computational observation, not a theorem.
\end{definition}

\begin{remark}
If the double transversality is structural, the
collision invariant has base-specific content that
cannot be recovered from the aggregate. The
convergence at each base would then be a local
phenomenon, depending on the digit geometry at that
base rather than on a universal property.
\end{remark}

% ============================================================
\section{Interpretation}

\begin{remark}
If the double transversality holds generally, the
convergence of $F^{\circ}_b(s)$ at a specific base $b$
is a local phenomenon, depending on the anti-correlation
at that base rather than on a universal property across
bases. This would separate the collision transform from
approaches that study families of $L$-functions through
averaged quantities: the base sum loses the structural
content that is present at each individual base.
\end{remark}

\begin{remark}
The vanishing of $F_R$ does not mean the base sum is
uninteresting. The neutrality theorem (the mod-$3$
splitting of the base sum) is a statement about $R(p)$
itself, not about $F_R(s)$. The double transversality
operates at the level of the prime harmonic sum, not
at the level of individual primes.
\end{remark}

% ============================================================
\section{Remarks}

\subsection*{What is proved, what is computed}

The higher-lag anti-correlation and the vanishing base
sum are computational observations. The anti-correlation
is tested at $3$ bases across $15$ lag values, from $3$
to $2{,}187$ odd characters. The base sum is computed
across $10$ bases with $148{,}000$ primes.

The general neutrality theorem and the scramble
proposition (from the preceding paper~\cite{paper19})
are proved. The double transversality is a name for
the observed structure, not a theorem.

\subsection*{What the higher lags reveal}

As the lag $\ell$ increases, the modulus
$m = b^{\ell+1}$ grows and the family of odd
$L$-functions expands. The anti-correlation persists
throughout the tested range, meaning the collision
invariant continues to avoid the zero-sensitive
directions in these computations even as the character
space grows. This is consistent with the avoidance
being structural rather than a small-modulus artifact,
but the paper does not prove that conclusion.

\subsection*{The arc}

The collision invariant began as a count of digit
coincidences~\cite{paperA}. The collision transform
revealed its character decomposition and proved
convergence at $s = 1$~\cite{paperB}. The general
neutrality theorem constrained the decomposition at
every prime scale~\cite{paper19}. The sign structure
and anti-correlation suggested that convergence below
$s = 1$ may be a collective phenomenon, not a test of
individual $L$-functions~\cite{paper19}.

The double transversality extends that picture. The
computations suggest the collision invariant's
relationship to $L$-function zeros is structured at
two levels: within each base (the spectrum is
anti-correlated with zero-sensitive characters) and
across bases (the individual signals cancel).
Whether this pattern reflects a structural property
of the digit function's geometry or an artifact of
finite computation remains open.

% ============================================================
\begin{thebibliography}{99}

\bibitem{davenport}
H.~Davenport,
\emph{Multiplicative Number Theory},
3rd~ed., Springer, 2000.

\bibitem{iwaniec}
H.~Iwaniec and E.~Kowalski,
\emph{Analytic Number Theory},
AMS Colloquium Publications, vol.~53, 2004.

\bibitem{paperA}
A.~S.~Petty,
The collision invariant,
preprint, 2026.

\bibitem{paperB}
A.~S.~Petty,
The collision transform,
preprint, 2026.

\bibitem{paper19}
A.~S.~Petty,
The general neutrality theorem,
research note, 2025.

\bibitem{nfield}
A.~S.~Petty,
\emph{nfield: A structural analysis engine for fractional
field invariants},
software.
\url{https://github.com/alexspetty/nfield}

\end{thebibliography}

\end{document}
